\documentclass{article}
\usepackage{lmodern}
\usepackage{amsmath}
\usepackage{listings}
\usepackage{fullpage}
\usepackage{parskip}
\usepackage{xcolor}
\lstset{
	basicstyle=\ttfamily,
	columns=fullflexible,
	frame=single,
	breaklines=true,
	postbreak=\mbox{\textcolor{red}{$\hookrightarrow$}\space},
}
\begin{document}
\title{Experiment `p\_6\_6\_alternating\_sum\_array' Results}
\author{\today}
\date{}
\maketitle
\textbf{Experiment outcome: }SUCCESS
\\\textbf{Bad responses: }0
\\\textbf{Responses containing}~\texttt{assume}~\textbf{: }0
\\\textbf{Responses containing}~\texttt{expect}~\textbf{: }0
\\\textbf{Resolution attempts: }1
\\\textbf{Hard fails (resolution): }0
\\\textbf{Soft fails (resolution): }0
\\\textbf{Verification attempts: }1
\section*{Problem Specification}
\textbf{Problem name: }p\_6\_6\_alternating\_sum\_array
\\\textbf{Natural language statement: }null
\\\textbf{Method signature: }p\_6\_6\_alternating\_sum\_array(inputs: array<int>) returns (alt\_sum: int)
\subsection*{Ensures}
\begin{itemize}
\item \texttt{alt\_sum == sum\_array(inputs[..], 1)}
\end{itemize}
\subsection*{Functional Code Given}
\begin{lstlisting}
function sum_array(inputs: seq<int>, i: int): int
{
  if |inputs| == 0 then 0
  else inputs[0] * i + sum_array(inputs[1..], i * -1)
}
\end{lstlisting}
\clearpage
\section*{GenAI interactions}
Below you will find all interactions between the `user' (program) and the `assistant' (OpenAI).
\subsection*{Program $\rightarrow$ GenAI}
\begin{lstlisting}
You are given the following task to perform in Dafny:

Write a method to satisfy the below requirements.

The signature should be:

p_6_6_alternating_sum_array(inputs: array<int>) returns (alt_sum: int)

The method should respect the following contract:

ensures alt_sum == sum_array(inputs[..], 1)

The contract uses the following dafny code:

function sum_array(inputs: seq<int>, i: int): int
{
  if |inputs| == 0 then 0
  else inputs[0] * i + sum_array(inputs[1..], i * -1)
}

These function(s) must not be used in your implementation of the method.

Produce and show only the Dafny body of this method, including the curly braces that surround it. Do not show the signature nor contract. You must not use 'assume' nor 'expect' anywhere in your code.
\end{lstlisting}
\subsection*{GenAI $\rightarrow$ Program}
\textbf{System fingerprint: }null
\\\textbf{ID: }chatcmpl-DnOMXYe3jvgpfazOq83ENUg27C0mG
\\\textbf{Model: }gpt-5.5-2026-04-23
\\\textbf{Created at: }1780664093
\\\textbf{Finish reason: }stop
\\\textbf{Usage: }410 tokens in, and 4540 tokens out
\begin{lstlisting}
{
  alt_sum := 0;
  var idx := 0;
  var sign := 1;

  assert inputs[0..] == inputs[..];
  assert alt_sum + sum_array(inputs[idx..], sign) == sum_array(inputs[..], 1);

  while idx < inputs.Length
    invariant 0 <= idx <= inputs.Length
    invariant alt_sum + sum_array(inputs[idx..], sign) == sum_array(inputs[..], 1)
    decreases inputs.Length - idx
  {
    var oldIdx := idx;
    var oldSign := sign;
    var rest := inputs[oldIdx..];

    assert 0 <= oldIdx < inputs.Length;
    assert |rest| > 0;
    assert rest[0] == inputs[oldIdx];
    assert rest[1..] == inputs[oldIdx + 1..];
    assert sum_array(rest, oldSign) == rest[0] * oldSign + sum_array(rest[1..], oldSign * -1);
    assert sum_array(inputs[oldIdx..], oldSign) == inputs[oldIdx] * oldSign + sum_array(inputs[oldIdx + 1..], oldSign * -1);

    alt_sum := alt_sum + inputs[oldIdx] * oldSign;
    sign := oldSign * -1;
    idx := oldIdx + 1;

    assert alt_sum + sum_array(inputs[idx..], sign) == sum_array(inputs[..], 1);
  }

  assert idx == inputs.Length;
  assert inputs[idx..] == [];
  assert sum_array(inputs[idx..], sign) == 0;
  assert alt_sum == sum_array(inputs[..], 1);
}
\end{lstlisting}
\section*{Final Program}
\begin{lstlisting}
function sum_array(inputs: seq<int>, i: int): int
{
  if |inputs| == 0 then 0
  else inputs[0] * i + sum_array(inputs[1..], i * -1)
}

method p_6_6_alternating_sum_array(inputs: array<int>) returns (alt_sum: int)
	ensures alt_sum == sum_array(inputs[..], 1)
{
  alt_sum := 0;
  var idx := 0;
  var sign := 1;

  assert inputs[0..] == inputs[..];
  assert alt_sum + sum_array(inputs[idx..], sign) == sum_array(inputs[..], 1);

  while idx < inputs.Length
    invariant 0 <= idx <= inputs.Length
    invariant alt_sum + sum_array(inputs[idx..], sign) == sum_array(inputs[..], 1)
    decreases inputs.Length - idx
  {
    var oldIdx := idx;
    var oldSign := sign;
    var rest := inputs[oldIdx..];

    assert 0 <= oldIdx < inputs.Length;
    assert |rest| > 0;
    assert rest[0] == inputs[oldIdx];
    assert rest[1..] == inputs[oldIdx + 1..];
    assert sum_array(rest, oldSign) == rest[0] * oldSign + sum_array(rest[1..], oldSign * -1);
    assert sum_array(inputs[oldIdx..], oldSign) == inputs[oldIdx] * oldSign + sum_array(inputs[oldIdx + 1..], oldSign * -1);

    alt_sum := alt_sum + inputs[oldIdx] * oldSign;
    sign := oldSign * -1;
    idx := oldIdx + 1;

    assert alt_sum + sum_array(inputs[idx..], sign) == sum_array(inputs[..], 1);
  }

  assert idx == inputs.Length;
  assert inputs[idx..] == [];
  assert sum_array(inputs[idx..], sign) == 0;
  assert alt_sum == sum_array(inputs[..], 1);
}
\end{lstlisting}
\section*{Total Token Usage}
\textbf{Input tokens: }410
\\\textbf{Output tokens: }4540
\\\textbf{Reasoning tokens: }4095
\\\textbf{Sum of `total tokens': }4950
\section*{Experiment Timings}
\textbf{Overall Experiment} started at 1780664092979, ended at 1780664207277, lasting 114298ms (114.30 seconds)
\\\textbf{Iteration \#1} started at 1780664092979, ended at 1780664207277, lasting 114298ms (114.30 seconds)
\end{document}
